\documentclass{article}
\usepackage[T1]{fontenc}
\usepackage{lmodern}
\usepackage{graphicx}
\usepackage{amsmath,amssymb}
\usepackage[a4paper,margin=2cm]{geometry}
\usepackage[absolute,overlay]{textpos}
\setlength{\TPHorizModule}{1mm}
\setlength{\TPVertModule}{1mm}
\usepackage[indonesian]{babel}
\usepackage{hyperref}
\setlength{\emergencystretch}{2em}
% unit: Halaman 6

\begin{document}

% === Halaman 6 (python/native) ===

• Lakukan proses yang sama untuk blok P2, P3, dst.

• Tentu saja algoritma enkripsi E dan D di atas sangat lemah karena mudah 

dipecahkan.

• Cipher blok modern membuat E dan D sangat rumit prosesnya sehingga sangat 

sukar dipecahkan oleh kriptanalis.

• Fungsi E dan D melibatkan sejumlah teknik seperti teknik substitusi, permutasi, 

pergeseran (shifting), kompresi, ekspansi, dsb.

• Setiap blok tidak dienkripsi satu kali, tetapi berulang kali untuk mendapatkan 

cipherteks yang kuat.

6

\end{document}
